
This paper tested whether category-specific temperature scaling could exploit observed calibration variation in LLMs. Experiments on TruthfulQA with Llama-2-7B confirmed that calibration error varies significantly across semantic categories (ANOVA p=0.00012, ECE range 0.099). However, this variation cannot be exploited for improved calibration: all categories require identical maximum temperature smoothing (T=10.0, CV=0), revealing uniform overconfidence that dominates category-specific effects.

This negative result constitutes the primary contribution. The paper provides the first systematic per-category calibration analysis on TruthfulQA, demonstrating that while category structure exists in calibration behavior, it does not translate to category-specific calibration needs. Future work should explore calibration methods beyond temperature scaling that can address the uniform severity of overconfidence across all categories, and should extend the temperature search range to determine whether cluster differentiation emerges at higher temperatures.
